\paperchapter{September 2018}{30 September 2018}{A}{50}{13th GSET}
\keytable{\kcell{1}{B} & \kcell{2}{A} & \kcell{3}{B} & \kcell{4}{C} & \kcell{5}{B} & \kcell{6}{B} & \kcell{7}{A} & \kcell{8}{D} & \kcell{9}{C} & \kcell{10}{B} \\
\kcell{11}{D} & \kcell{12}{B} & \kcell{13}{D} & \kcell{14}{C} & \kcell{15}{C} & \kcell{16}{A} & \kcell{17}{B} & \kcell{18}{C} & \kcell{19}{D} & \kcell{20}{B} \\
\kcell{21}{B} & \kcell{22}{A} & \kcell{23}{C} & \kcell{24}{B} & \kcell{25}{A} & \kcell{26}{A} & \kcell{27}{D} & \kcell{28}{C} & \kcell{29}{C} & \kcell{30}{D} \\
\kcell{31}{D} & \kcell{32}{C} & \kcell{33}{B} & \kcell{34}{B} & \kcell{35}{C} & \kcell{36}{C} & \kcell{37}{A} & \kcell{38}{C} & \kcell{39}{A} & \kcell{40}{D} \\
\kcell{41}{A} & \kcell{42}{C} & \kcell{43}{A} & \kcell{44}{B} & \kcell{45}{A} & \kcell{46}{C} & \kcell{47}{B} & \kcell{48}{C} & \kcell{49}{C} & \kcell{50}{D}}
\begin{question}{1}
{Name the philosopher who proposed the technique of questioning:\par}
\optsfour{John Dewey}{Socrates}{Comte}{Froebel}
\answer{B}{Socrates}
\begin{explain}
\textbf{Socrates} (469--399 BCE) taught by systematic questioning. The \emph{Socratic method} (maieutics, the ``midwife'' method) leads learners, through a chain of questions, to discover contradictions in their views and arrive at knowledge themselves.
\why{Dewey gave the problem/project method, Comte founded positivism, and Froebel founded the kindergarten (play-way method).}
\end{explain}
\end{question}

\begin{question}{2}
{Which of the following skills does not require high intellectual ability?\par}
\optstwo{Technical skills}{Mechanical skills}{Teaching skills}{Management skills}
\answer{A}{Technical skills}
\begin{explain}
According to the official key, \textbf{technical skills} (routine operational or procedural skills) require the least intellectual ability. Teaching and management skills demand high-level conceptual and interpersonal thinking.
\end{explain}
\end{question}

\begin{question}{3}
{Primary education in the mother tongue has been safeguarded in Article :\par}
\optsfour{29}{350}{220}{351}
\answer{B}{350}
\begin{explain}
\textbf{Article 350A} of the Constitution directs every State to provide adequate facilities for instruction in the \textbf{mother tongue at the primary stage} to children of linguistic minority groups.
\why{Article 29 protects the culture and language of minorities in general; Article 351 concerns the development of Hindi.}
\end{explain}
\end{question}

\begin{question}{4}
{What is the full form of CAI ?\par}
\optstwo{Computer Associated Instruction}{Computer Associated Information}{Computer Assisted Instruction}{Computer Assisted Information}
\answer{C}{Computer Assisted Instruction}
\begin{explain}
\textbf{CAI} is \textbf{Computer Assisted Instruction}: using computers to present instructional material, drill and practice, tutorials and simulations, with immediate feedback. It grew out of Skinner's programmed learning.
\end{explain}
\end{question}

\begin{question}{5}
{For maintaining discipline in the class:\par}
\optsone{Students should be severely punished for acts of indiscipline}{Teacher should try to make his teaching interesting}{Non-attending students should be turned out of the class}{Parents of the indisciplined students should be informed}
\answer{B}{Teacher should try to make his teaching interesting}
\begin{explain}
The best preventive discipline is \textbf{interesting, well-planned teaching}. When learners are engaged, indiscipline rarely arises.
\why{Severe punishment, expulsion from class or complaining to parents are negative measures that do not address the cause.}
\end{explain}
\end{question}

\begin{question}{6}
{If the sample drawn does not specify any condition about the parameter of the population, it is called :\par}
\optstwo{Selected statistics}{Distribution free statistics}{Parametric statistics}{Inferential statistics}
\answer{B}{Distribution free statistics}
\begin{explain}
Statistical tests that make \textbf{no assumptions about population parameters} (such as normality or equal variance) are called \textbf{non-parametric} or \textbf{distribution-free} statistics. Examples are chi-square, Mann--Whitney U, Wilcoxon, Kruskal--Wallis and Spearman's rho.
\end{explain}
\end{question}

\begin{question}{7}
{Validity of an experimental design refers to:\par}
\optsone{Maximise treatment variance and minimise error variance}{Minimise both treatment variance and error variance}{Testing the difference between treatment variance and error variance}{Minimise treatment variance and maximise error variance}
\answer{A}{Maximise treatment variance and minimise error variance}
\begin{explain}
Kerlinger's \textbf{MAX-MIN-CON} principle for a good experimental design: \textbf{MAX}imise systematic (treatment) variance, \textbf{MIN}imise error variance and \textbf{CON}trol extraneous variance. A design that does this yields valid conclusions.
\end{explain}
\end{question}

\begin{question}{8}
{Which one of the following is not considered as plagiarism?\par}
\optsone{Make use of the work of another and misrepresent it as your own}{Drawing content from the work of another without acknowledging the source}{Paraphrasing too closely to the original text}{Drawing content from another work and adapting it with due acknowledgement}
\answer{D}{Drawing content from another work and adapting it with due acknowledgement}
\begin{explain}
Using someone else's work \textbf{with proper acknowledgement} (citation) is legitimate scholarship, not plagiarism.
\why{Passing off others' work as one's own, using content without acknowledgement, and paraphrasing too closely are all forms of plagiarism.}
\end{explain}
\end{question}

\begin{question}{9}
{Bibliography refers to:\par}
\optsone{A true story written about someone}{A record of events presented in a report}{A list of sources used in a report and where they can be found}{List of abbreviations used in the report}
\answer{C}{A list of sources used in a report and where they can be found}
\begin{explain}
A \textbf{bibliography} is a list of all the \emph{sources} (books, articles, websites) consulted or cited, with enough detail (author, title, publisher, year, pages, URL) for readers to locate them.
\why{A true story about someone is a \emph{biography}; a list of abbreviations is a glossary or abbreviation list.}
\end{explain}
\end{question}

\begin{question}{10}
{Snowball sampling technique falls under the category of:\par}
\optstwo{Probability sampling}{Non-probability sampling}{Random sampling}{Quota sampling}
\answer{B}{Non-probability sampling}
\begin{explain}
\textbf{Snowball sampling} is a \textbf{non-probability} technique: existing participants recruit further participants from among their acquaintances. It is useful for hidden populations.
\end{explain}
\end{question}

\begin{context}
\textbf{Question Nos.~11--15 should be attempted from the following paragraph.} Below is given a passage followed by possible inferences which can be drawn from the facts stated in the passage. You have to examine each inference separately in the context of the passage and decide upon its degree of truth or falsity.\par\smallskip
\textbf{Passage:} ``The fifth conference of Labour Ministers of Non-Aligned and other developing countries which concluded in New Delhi on January 23, was not just another conference held under the auspices of International Labour Organization to discuss labour welfare. For one it was in the run-up to the world summit on social development that the conference was held, and for another, it was under the shadow of the demand by the industrialised countries to link labour standards to International trade through a social cause. The conference covered a large part of the agenda to be discussed at the summit to be held at Copenhagen in March, i.e.\ employment creation, poverty reduction and eradication and social integration.''
\end{context}
\begin{question}{11}
{The fourth conference of Labour Ministers was not held at New Delhi:\par}
\optstwo{The inference is probably true}{The inference is probably false}{The inference is definitely false}{The data are inadequate}
\answer{D}{The data are inadequate}
\begin{explain}
The passage tells us only where the \emph{fifth} conference was held (New Delhi). It says nothing about the venue of the fourth conference, so the \textbf{data are inadequate}.
\end{explain}
\end{question}

\begin{question}{12}
{blank{} organization took an active part in a conference of Labour Ministers for the first time.\par}
\optstwo{International Monetary Fund}{International Labour Organization}{World Health Organization}{UNICEF}
\answer{B}{International Labour Organization}
\begin{explain}
The conference was held ``under the auspices of the \textbf{International Labour Organization}''. The ILO is the organisation concerned with labour; the IMF, WHO and UNICEF are not mentioned. The key accepts (B).
\end{explain}
\end{question}

\begin{question}{13}
{Linking Labour standards to International trade would benefit the blank{} countries.\par}
\optsfour{Advanced}{Developing}{Developed}{Industrial}
\answer{D}{Industrial}
\begin{explain}
The demand to link labour standards to trade came from the \textbf{industrialised countries}. It would help them by making the products of low-wage developing countries less competitive, so the linkage would benefit the industrial countries.
\end{explain}
\end{question}

\begin{question}{14}
{Among the following topics which of them is \emph{not} discussed at the Copenhagen Conference?\par}
\optstwo{Employment}{Poverty}{Technology}{Social Integration}
\answer{C}{Technology}
\begin{explain}
The Copenhagen agenda was ``employment creation, poverty reduction and eradication and social integration''. \textbf{Technology} is not mentioned.
\end{explain}
\end{question}

\begin{question}{15}
{Which of the following is the reason for organizing the fifth conference of Labour Ministers?\par}
\optstwo{Compulsion by ILO}{To satisfy the government}{To run-up to the world summit}{Causal meeting of Ministers}
\answer{C}{To run-up to the world summit}
\begin{explain}
``\ldots it was \textbf{in the run-up to the world summit} on social development that the conference was held.''
\end{explain}
\end{question}

\begin{question}{16}
{Message is any signal that triggers the response of a\par}
\optsfour{Receiver}{Sender}{Channel}{Feedback}
\answer{A}{Receiver}
\begin{explain}
A message is any signal (verbal or non-verbal) sent by the source that stimulates a response in the \textbf{receiver}.
\end{explain}
\end{question}

\begin{question}{17}
{Our dress code is an example of \ldots{} communication.\par}
\optsfour{Verbal}{Non-verbal}{Written}{Spoken}
\answer{B}{Non-verbal}
\begin{explain}
Dress, appearance and personal artefacts communicate without words, so they are \textbf{non-verbal communication} (artifactics).
\end{explain}
\end{question}

\begin{question}{18}
{\ldots{} is the process by which the receiver interprets the symbols used by the source of the message.\par}
\optsfour{Listing}{Listening}{Decoding}{Encoding}
\answer{C}{Decoding}
\begin{explain}
\textbf{Decoding} is the receiver's interpretation of the symbols (words, signs, gestures) used by the sender, i.e.\ converting the message back into meaning.
\end{explain}
\end{question}

\begin{question}{19}
{Identify the correct sequence in the communication process:\par}
\optstwo{Sender - Channel - Message - Receiver}{Sender - Receiver - Message - Channel}{Sender - Receiver - Channel - Message}{Sender - Message - Channel - Receiver}
\answer{D}{Sender - Message - Channel - Receiver}
\begin{explain}
The communication process runs \textbf{Sender $\rightarrow$ Message $\rightarrow$ Channel $\rightarrow$ Receiver} (Berlo's SMCR model), with feedback returning to the sender.
\end{explain}
\end{question}

\begin{question}{20}
{The study of body language is called\par}
\optsfour{Proxemics}{Kinesics}{Phonetics}{Chronemics}
\answer{B}{Kinesics}
\begin{explain}
\textbf{Kinesics} (Ray Birdwhistell, 1952) is the study of body language: gestures, posture, facial expressions and eye movements.
\why{Proxemics = space; Phonetics = speech sounds; Chronemics = time.}
\end{explain}
\end{question}

\begin{question}{21}
{Find out missing number in the following:\par}
\qfigure{\petal{5}{3}{6}{8}{63}\quad\petal{2}{7}{3}{9}{41}\quad\petal{6}{7}{8}{5}{?}}
\optsfour{26}{82}{83}{86}
\answer{B}{82}
\begin{explain}
Centre $=$ (product of the top pair) $+$ (product of the bottom pair):
\[5\times3+6\times8=15+48=63,\qquad 2\times7+3\times9=14+27=41\]
\[6\times7+8\times5=42+40=\mathbf{82}\]
\end{explain}
\end{question}

\begin{question}{22}
{If P means x, R means +, T means $\div$ and S means -, then 18T3P9S8R6 = ?\par}
\optsfour{52}{2/3}{46}{58}
\answer{A}{52}
\begin{explain}
Replace the letters: $18\div3\times9-8+6$.
By BODMAS: $18\div3=6$; $6\times9=54$; $54-8+6=\mathbf{52}$.
\end{explain}
\end{question}

\begin{question}{23}
{Young is related to Old, in the same way Wide is related to\par}
\optsfour{Short}{Big}{Narrow}{Long}
\answer{C}{Narrow}
\begin{explain}
Young and old are \textbf{antonyms}; the antonym of wide is \textbf{narrow}.
\end{explain}
\end{question}

\begin{question}{24}
{If SUMMER is coded as RUNNER, the code for WINTER will be:\par}
\optsfour{SUITER}{VIOUER}{WALKER}{SUFFER}
\answer{B}{VIOUER}
\begin{explain}
SUMMER $\rightarrow$ RUNNER: S$\rightarrow$R ($-1$), U$\rightarrow$U, M$\rightarrow$N ($+1$), M$\rightarrow$N ($+1$), E$\rightarrow$E, R$\rightarrow$R.
Applying the same shifts to WINTER: W$-1=$V, I, N$+1=$O, T$+1=$U, E, R, giving \textbf{VIOUER}.
\end{explain}
\end{question}

\begin{question}{25}
{Find out the missing term in the following series : 0 7 26 63 124 \ldots{}\par}
\optsfour{215}{216}{248}{157}
\answer{A}{215}
\begin{explain}
Each term is $n^3-1$: $1-1=0$, $8-1=7$, $27-1=26$, $64-1=63$, $125-1=124$, so next $216-1=\mathbf{215}$.
\end{explain}
\end{question}

\begin{question}{26}
{In the following question choose the group of words that shows the same relationship as given at the top of the question : Play : Win : Lose\par}
\optstwo{Accident : Death : Survive}{Examination : Success : Determination}{Read : Book : Magazine}{Music : Dance : Art}
\answer{A}{Accident : Death : Survive}
\begin{explain}
Play has two opposite possible outcomes, win or lose. Likewise an \textbf{accident} may end in \textbf{death} or the victim may \textbf{survive}.
\why{In (B), determination is not an outcome of an examination; (C) and (D) are not outcome pairs.}
\end{explain}
\end{question}

\begin{question}{27}
{What is the missing number in the following sequence of numbers ? 84/14,12 81/18,9 88/?,11\par}
\optsfour{13}{14}{15}{16}
\answer{D}{16}
\begin{explain}
The number above each fraction is half the product of the two numbers below it:
\[\frac{14\times12}{2}=84,\qquad\frac{18\times9}{2}=81,\qquad\frac{?\times11}{2}=88\Rightarrow ?=\mathbf{16}\]
\end{explain}
\end{question}

\begin{question}{28}
{In a group of persons travelling in a bus, 6 persons can speak Malayalam, 15 can speak Hindi and 6 can speak Gujarati. In that group, none can speak any other language. If 2 persons in the group can speak two languages and one person can speak all the three languages, then how many persons are there in the group?\par}
\optsfour{21}{22}{23}{24}
\answer{C}{23}
\begin{explain}
Adding the language groups ($6+15+6=27$) counts a person who speaks exactly two languages twice and the person who speaks all three three times.
\[\text{Persons}=27-(\text{exactly two})-2(\text{all three})=27-2-2\times1=\mathbf{23}\]
\end{explain}
\end{question}

\begin{question}{29}
{In this question, two statements followed by four conclusions are given...\newline Statements:\newline (1) Some doors are mangoes.\newline (2) All mangoes are bananas.\newline Conclusions :\newline (I) All bananas are mangoes.\newline (II) All doors are bananas.\newline (III) Some doors are bananas.\newline (IV) Some mangoes are doors.\newline Which of the given conclusions logically follows...?\par}
\optstwo{Only I and II follows}{Only II and III follows}{Only III and IV follows}{Only II and IV follows}
\answer{C}{Only III and IV follows}
\begin{explain}
\begin{center}
\begin{tikzpicture}[scale=0.8,font=\scriptsize]
\draw (0,0) circle(1.7); \node at (0,1.35){Bananas};
\draw[fill=brand!10] (0.55,-0.25) circle(0.8); \node at (0.8,-0.25){Mangoes};
\draw (-1.5,-0.3) ellipse(1.5 and 0.6); \node at (-2.3,-0.3){Doors};
\end{tikzpicture}
\end{center}
Some doors are mangoes, and all mangoes are bananas, so \textbf{some doors are bananas} (III follows). ``Some doors are mangoes'' converts to \textbf{some mangoes are doors} (IV follows). ``All bananas are mangoes'' (I) and ``All doors are bananas'' (II) do not follow. The answer is \textbf{only III and IV}.
\end{explain}
\end{question}

\begin{question}{30}
{The pie-diagram below shows various expenditures incurred in publishing a book:\par}
\qfigure{\begin{tikzpicture}\pie[radius=1.3,text=pin,font=\scriptsize,color={brand!10,brand!35,brand!20,brand!50,brand!5,brand!25}]{15/Royalty,20/Printing,25/Paper,30/Binding,10/Promotion,10/Transportation}\end{tikzpicture}}
If for an edition of the book, the cost of paper is Rs.~20,000, then what is the promotion cost for this edition?
\optsfour{Rs. 5,000}{Rs. 6,000}{Rs. 7,000}{Rs. 8,000}
\answer{D}{Rs. 8,000}
\begin{explain}
Paper $=25\%$ corresponds to Rs.~20,000, so $1\%=$ Rs.~800.
Promotion $=10\%=10\times800=$ \textbf{Rs.~8,000}.
\end{explain}
\end{question}

\begin{question}{31}
{If ${}^{2n}C_{n-1} : {}^{2n-1}C_{n} = 5:3$, then what is the value of $n$?\par}
\optsfour{2}{3}{4}{5}
\answer{D}{5}
\begin{explain}
\[\frac{{}^{2n}C_{n-1}}{{}^{2n-1}C_n}=\frac{\frac{(2n)!}{(n-1)!\,(n+1)!}}{\frac{(2n-1)!}{n!\,(n-1)!}}=\frac{2n\cdot n!}{(n+1)!}=\frac{2n}{n+1}\]
\[\frac{2n}{n+1}=\frac53\Rightarrow6n=5n+5\Rightarrow n=\mathbf{5}\]
\end{explain}
\end{question}

\begin{question}{32}
{The average of a set of 7 observations is 36 and that of another set of 13 observations is 46. What is the combined average of the set ?\par}
\optsfour{40.5}{41.5}{42.5}{43.5}
\answer{C}{42.5}
\begin{explain}
\[\bar x=\frac{7\times36+13\times46}{7+13}=\frac{252+598}{20}=\frac{850}{20}=\mathbf{42.5}\]
\end{explain}
\end{question}

\begin{question}{33}
{The following circuit operates only if there is a path of functional devices from left to right. The proportion that each device functions is shown on the graph. Assume that devices fail independently. What is the proportion that the circuit operates?\par}
\qfigure{\begin{tikzpicture}[font=\scriptsize,box/.style={draw,minimum width=9mm,minimum height=5mm}]
\node (A) at (0,0) {A}; \node[box] (a) at (1.5,0.8) {0.90}; \node[box] (b) at (1.5,0) {0.79}; \node[box] (c) at (1.5,-0.8) {0.95};
\node[box] (d) at (3.6,0.45) {0.95}; \node[box] (e) at (3.6,-0.45) {0.98}; \node[box] (f) at (5.6,0) {0.99}; \node (B) at (6.8,0) {B};
\draw (A) -- (0.7,0); \draw (0.7,0.8) -- (0.7,-0.8); \foreach \n in {a,b,c} {\draw (0.7,0 |- \n.west) -- (\n.west); \draw (\n.east) -- (2.3,0 |- \n.east);}
\draw (2.3,0.8) -- (2.3,-0.8); \draw (2.3,0) -- (2.8,0); \draw (2.8,0.45) -- (2.8,-0.45);
\foreach \n in {d,e} {\draw (2.8,0 |- \n.west) -- (\n.west); \draw (\n.east) -- (4.4,0 |- \n.east);}
\draw (4.4,0.45) -- (4.4,-0.45); \draw (4.4,0) -- (f.west); \draw[-Stealth] (f.east) -- (B);
\end{tikzpicture}}
\optsfour{0.9725}{0.9879}{0.9915}{0.9998}
\answer{B}{0.9879}
\begin{explain}
A parallel block works unless \emph{all} its devices fail; blocks in series must \emph{all} work.
\[P_1=1-(0.10)(0.21)(0.05)=1-0.00105=0.99895\]
\[P_2=1-(0.05)(0.02)=0.999,\qquad P_3=0.99\]
\[P=0.99895\times0.999\times0.99\approx\mathbf{0.9879}\]
\end{explain}
\end{question}

\begin{question}{34}
{The value of $\dfrac{6^{2/3}\times\sqrt[3]{6^7}}{\sqrt[3]{6^6}}$ is:\par}
\optsfour{3}{6}{9}{12}
\answer{B}{6}
\begin{explain}
\[\frac{6^{2/3}\times6^{7/3}}{6^{6/3}}=6^{\frac23+\frac73-2}=6^{3-2}=\mathbf{6}\]
\end{explain}
\end{question}

\begin{question}{35}
{A bag contains 2 red, 3 green and 2 blue balls. Two balls are drawn at random. What is the probability that none of the balls drawn is blue ?\par}
\optsfour{$\dfrac27$}{$\dfrac{11}{21}$}{$\dfrac{10}{21}$}{$\dfrac57$}
\answer{C}{$\dfrac{10}{21}$}
\begin{explain}
Total balls $=7$, of which non-blue $=5$ (2 red $+$ 3 green).
\[P(\text{no blue})=\frac{{}^5C_2}{{}^7C_2}=\frac{10}{21}\]
\end{explain}
\end{question}

\begin{question}{36}
{CD-ROM stands for:\par}
\optstwo{Compactable Read Only Memory}{Compactable Data Read Only Memory}{Compact Disk Read Only Memory}{Common Disk Read Only Memory}
\answer{C}{Compact Disk Read Only Memory}
\begin{explain}
\textbf{CD-ROM} is \textbf{Compact Disc Read-Only Memory}, an optical disc holding about 700 MB that can be read but not rewritten.
\end{explain}
\end{question}

\begin{question}{37}
{WAN stands for:\par}
\optstwo{Wide Area Network}{Wireless Access Network}{Wireless Area Network}{Wide Access Network}
\answer{A}{Wide Area Network}
\begin{explain}
\textbf{WAN} is \textbf{Wide Area Network}, a network spanning large geographical areas (cities or countries). The Internet is the largest WAN.
\tip{PAN $<$ LAN $<$ MAN $<$ WAN, in increasing geographical coverage.}
\end{explain}
\end{question}

\begin{question}{38}
{A computer cannot 'boot', if it does not have the :\par}
\optstwo{Compiler}{Loader}{Operating System}{Keypad}
\answer{C}{Operating System}
\begin{explain}
\textbf{Booting} means loading the \textbf{operating system} into main memory when the computer starts (the BIOS/UEFI runs POST and then the boot loader). Without an OS the computer cannot boot.
\end{explain}
\end{question}

\begin{question}{39}
{Human beings are referred to as Homo sapiens. Which device is called Sillico sapiens ?\par}
\optsfour{Computer}{Hardware}{Robot}{Drone}
\answer{A}{Computer}
\begin{explain}
The \textbf{computer} is popularly called \emph{Silico sapiens}, a ``wise'' (thinking) machine built on silicon chips, just as humans are \emph{Homo sapiens}.
\end{explain}
\end{question}

\begin{question}{40}
{1 GB is equal to\par}
\optsfour{1250 megabytes}{1296 megabytes}{1150 megabytes}{1024 megabytes}
\answer{D}{1024 megabytes}
\begin{explain}
In the binary system, $1\ \text{GB}=2^{10}\ \text{MB}=\mathbf{1024\ MB}$.
\end{explain}
\end{question}

\begin{question}{41}
{Which is India's first bullet train project?\par}
\optstwo{Mumbai-Ahmedabad}{Delhi-Chandigarh}{Delhi-Jaipur}{Chennai-Kolkata}
\answer{A}{Mumbai-Ahmedabad}
\begin{explain}
India's first high-speed rail (bullet train) project is the \textbf{Mumbai--Ahmedabad High Speed Rail} corridor (about 508 km), built with Japanese Shinkansen technology. Its foundation stone was laid in September 2017.
\end{explain}
\end{question}

\begin{question}{42}
{Supreme Court of India has ordered the transition of vehicles to which BS level?\par}
\optsfour{BS-V}{BS-III}{BS-VI}{BS-IV}
\answer{C}{BS-VI}
\begin{explain}
In 2016 the Government decided to leapfrog from BS-IV directly to \textbf{BS-VI} emission norms. In October 2018 the Supreme Court ordered that no BS-IV vehicle be sold or registered from 1 April 2020, so the transition was to \textbf{BS-VI}.
\end{explain}
\end{question}

\begin{question}{43}
{Scientists have developed a new battery material based in \ldots{} with charging speed 5 times faster than lithium-ion batteries.\par}
\optsfour{Graphene}{Carbon}{Nitrate}{Chromium}
\answer{A}{Graphene}
\begin{explain}
Researchers (notably at the Samsung Advanced Institute of Technology) developed a \textbf{graphene}-based battery material (``graphene balls'') that allows charging about \textbf{five times faster} than conventional lithium-ion batteries.
\end{explain}
\end{question}

\begin{question}{44}
{Which is regarded as one of the world's hottest hotspots of biological diversity ?\par}
\optsfour{Rann of Kutch}{Western Ghats}{Eastern Ghats}{Thar Deserts}
\answer{B}{Western Ghats}
\begin{explain}
The \textbf{Western Ghats} (with Sri Lanka) form one of the world's 36 biodiversity hotspots and are regarded as one of the ``hottest'', with very high endemism. They are a UNESCO World Heritage Site (2012).
\tip{India's four hotspots: Himalaya, Indo-Burma, Western Ghats--Sri Lanka, and Sundaland (Nicobar Islands).}
\end{explain}
\end{question}

\begin{question}{45}
{Which Indian astrophysicist and Nobel laureate predicted rapidly rotating stars emitting polarised light ?\par}
\optstwo{S. Chandrasekhar}{C.V. Raman}{S. Ramanujam}{Venkat Raman Ramakrishnan}
\answer{A}{S. Chandrasekhar}
\begin{explain}
\textbf{S.~Chandrasekhar} (Nobel Prize in Physics, 1983) predicted that rapidly rotating stars would emit \textbf{polarised light}. Astronomers confirmed the prediction in 2017, about 70 years later, observing the star Regulus.
\end{explain}
\end{question}

\begin{question}{46}
{What is the full form of NIRF a methodology to rank higher education institutions in India?\par}
\optsone{National Institute of Research Framework}{National and International Ranking Framework}{National Institutional Ranking Framework}{National Inventory of Research Funding}
\answer{C}{National Institutional Ranking Framework}
\begin{explain}
\textbf{NIRF} is the \textbf{National Institutional Ranking Framework} (MHRD, 2015), which ranks HEIs on five parameters: TLR, RPP, GO, OI and Perception.
\end{explain}
\end{question}

\begin{question}{47}
{'SWAYAM' is a portal initiated by:\par}
\optstwo{NITI Aayog}{MHRD}{Ministry of I\&B}{Ministry of Railways}
\answer{B}{MHRD}
\begin{explain}
\textbf{SWAYAM} was initiated by the \textbf{Ministry of Human Resource Development (MHRD)}, now the Ministry of Education, and launched on 9 July 2017.
\end{explain}
\end{question}

\begin{question}{48}
{'AICTE' is responsible for the planning, coordination and management of:\par}
\optstwo{Teacher Education}{Distance Education}{Technical and Management Education}{School Education}
\answer{C}{Technical and Management Education}
\begin{explain}
The \textbf{AICTE} plans, coordinates and manages \textbf{technical and management education} (engineering, technology, management, pharmacy, architecture, hotel management and others).
\end{explain}
\end{question}

\begin{question}{49}
{In who's name the National Mission on Teachers and Teaching has been launched in India?\par}
\optstwo{Pt. Dindayal Upadhyaya}{Pt. Jawaharlal Nehru}{Pt. Madan Mohan Malaviya}{Pt. Ravishankar}
\answer{C}{Pt. Madan Mohan Malaviya}
\begin{explain}
The \textbf{Pandit Madan Mohan Malaviya National Mission on Teachers and Teaching (PMMMNMTT)} was launched in December 2014 to strengthen teacher education and faculty development at all levels.
\end{explain}
\end{question}

\begin{question}{50}
{In which state the first children's university was established in India?\par}
\optsfour{Maharashtra}{Rajasthan}{Madhya Pradesh}{Gujarat}
\answer{D}{Gujarat}
\begin{explain}
India's first \textbf{Children's University} was established at \textbf{Gandhinagar, Gujarat} in 2009 under the Children's University Act. It carries out research on child development and education.
\end{explain}
\end{question}
